This report presents the \emph{LLM-as-a-Judge Reliability Lab}, a reproducible study of whether large language models can act as dependable pairwise evaluators of generated responses. The central premise is deliberately modest: an evaluator verdict is an observation produced by a model, a prompt, a presentation order, and a protocol. It is not a ground-truth statement. The study therefore measures several distinct reliability properties instead of asking whether one judge is ``best.''

The final analysis uses a source-text-corrected controlled dataset with 1,568 records. It represents 80 distinct MT-Bench questions (upstream question IDs 81--160) as 160 turn-level prompts. Human preference reference labels are retained as the comparison target. The source-corrected evidence uses complete-case estimators, fixed seeds, 10,000 bootstrap iterations, no imputation, and provenance-preserving source precedence. Fourteen terminal Claude observations remain unavailable and are transparently excluded where a complete case is required.

Agreement with human preference reference labels was 59.20\% (457/772; 95\% CI 55.70--62.69\%; $\kappa=0.3230$). Repeated evaluations under fixed conditions were highly consistent (96.53\%, 7,119/7,375). In contrast, answer-order presentation mattered: 17.49\% of decisive paired comparisons changed their mapped outcome after an order swap (96/549). The narrowly defined redundant-text and presentation-format treatments produced no stable variant wins in their eligible pairs (0/682 and 0/716 respectively); these results are not general claims about verbosity or formatting. A counterbalanced source-family experiment estimated a 50.94\% stable same-family preference rate (81/159; 95\% CI 42.77--58.49\%) with strong judge-level heterogeneity.

Two mitigation families were evaluated as complementary, not head-to-head, strategies. The primary DUAL\_SWAP within-judge presentation-consistency filter had a +0.67 percentage-point agreement change (95\% CI -0.17 to +1.68) on its retained matched population, while reducing coverage by 21.03 points. The secondary strict four-judge consensus protocol improved agreement by +5.02 points (95\% CI +4.16 to +5.89) against its own equal-weight individual-judge comparator, retaining 69.83\% of planned pairs. The report documents the methods, findings, limitations, software architecture, and an evidence-led reproducibility workflow.

\textbf{Keywords:} LLM-as-a-Judge; reliability; human preference reference labels; position sensitivity; reproducibility; consensus; controlled evaluation.
